\documentclass[10pt]{article}
\usepackage[margin=0.85in]{geometry}
\usepackage{booktabs}
\usepackage{amsmath}
\usepackage{siunitx}
\usepackage{hyperref}
\usepackage{caption}
\usepackage{enumitem}
\setlist{nosep,leftmargin=*}
\setlength{\abovecaptionskip}{4pt}
\setlength{\belowcaptionskip}{2pt}

\title{Quantum Simulation of Photo-Induced Charge Dynamics:\\From Ethylene to Photosynthesis}
\author{MOSHARRAFA Team \\ Alexandria Quantum Hackathon 2026 --- Track 1}
\date{September 8, 2026}

\begin{document}
\maketitle

\begin{abstract}
We implement a VQE+QSE pipeline using the quantum equation-of-motion
(qEOM) method as the numerically robust realization of QSE to compute the
ground and low-lying excited states of ethylene (C\textsubscript{2}H\textsubscript{4})
in a (2$e$,2$o$) active space at the STO-3G level, benchmark it against
classical CASCI, and evaluate hardware noise and error mitigation on both
simulator and real IBM Quantum hardware. VQE agrees with CASCI to within
0.02~mHa; the triplet and singlet ($\pi\to\pi^*$) excited states are
correctly identified; and zero-noise extrapolation with readout mitigation
reduces the real-hardware ground-state error by a factor of $\sim$2.8. We
close with a proposed extension toward electron-transfer rates in a
chlorophyll-inspired system.
\end{abstract}

\section{Introduction}
Photo-induced electronic excitations underlie vision, photosynthesis, and
photovoltaic energy conversion, and are classically expensive to simulate as
molecules grow, motivating near-term quantum approaches such as VQE combined
with QSE, realized here via the qEOM method. We use ethylene's
well-characterized $\pi\to\pi^*$ excitation as a testbed for a pipeline
extensible toward larger, biologically relevant chromophores.

\section{Methodology}

\subsection{Molecular Setup and Active Space}
Ethylene was constructed at its equilibrium planar geometry
(C--C = 1.34~\AA, C--H = 1.087~\AA, H--C--H = 117.3\textdegree, verified
numerically against literature values) and at a nominally twisted geometry
intended to represent 90\textdegree\ rotation about the C--C axis. The active
space was restricted to 2 electrons in 2 spatial orbitals (the
$\pi$/$\pi^*$ pair) using \texttt{ActiveSpaceTransformer}, mapped to 4 qubits
via the Jordan--Wigner transformation.

\subsection{Ground-State VQE}
Two ansatze were implemented: UCCSD (3 parameters) and a hardware-efficient
\texttt{EfficientSU2} ansatz (24 parameters, linear entanglement, 2
repetitions), both optimized with COBYLA against a noiseless statevector
estimator.

\subsection{From QSE to QEOM}
The challenge statement specifies quantum subspace expansion (QSE): diagonalize
$H$ in the subspace spanned by excitation operators applied to the VQE ground
state, $\langle\psi_0|O_i^\dagger H O_j|\psi_0\rangle$. We implement this via
Qiskit Nature's \texttt{QEOM} (quantum equation-of-motion) solver rather than
a bare QSE diagonalization for two reasons. First, qEOM is the numerically
robust formulation of the same subspace idea: it builds the generalized
eigenvalue problem from double-commutator matrices ($M,Q,V,W$), which are
Hermitian by construction and better-conditioned under sampling noise than
directly diagonalizing $\langle O_i^\dagger H O_j\rangle$, which can become
ill-conditioned or lose Hermiticity from linear dependence and finite-shot
error. Second, qEOM is the version natively implemented and tested in Qiskit
Nature, avoiding a from-scratch, less-validated custom QSE implementation
under hackathon time constraints. The two methods are mathematically the same
subspace-expansion idea and are used interchangeably in the near-term quantum
chemistry literature; single and double excitations (``sd'') were used to
build the subspace on top of the converged VQE ground state.

\subsection{Classical Benchmarking}
CASCI(2,2)/STO-3G calculations (PySCF) were performed at both geometries as
an independent classical cross-check of the VQE and QEOM results.

\subsection{Noise-Aware Benchmarking and Error Mitigation}
Two noise studies were carried out: (i) the converged UCCSD ground state was
evaluated on real IBM Quantum hardware (\texttt{ibm\_marrakesh}) at
Qiskit Runtime resilience levels 0 (raw), 1 (TREX readout mitigation), and 2
(zero-noise extrapolation + TREX + gate twirling); (ii) the QEOM excited-state
calculation was repeated locally under a realistic device noise model
(derived from a \texttt{FakeSherbrooke} hardware snapshot) using Qiskit Aer's
\texttt{EstimatorV2}, isolating the effect of noise on the excitation-operator
measurements while holding the (noiselessly-optimized) ground state fixed.

\section{Results}

\subsection{Ground-State Energies}
Table~\ref{tab:gs} summarizes the ground-state results at the planar geometry.
Both ansatze agree with each other and with the classical CASCI(2,2)
reference to within 0.02~mHa, confirming correct implementation: with only 2
active electrons, UCCSD is variationally exact, and CASCI is the exact
diagonalization in the same active space.

\begin{table}[h]
\centering
\caption{Ground-state energy, circuit resources, and classical cross-check
(planar geometry).}
\label{tab:gs}
\begin{tabular}{lccc}
\toprule
Method & Energy (Ha) & Qubits & Depth (parameters) \\
\midrule
UCCSD (VQE) & $-77.116628$ & 4 & 112 (3) \\
Hardware-efficient (VQE)\footnotemark & $-77.116184$ & 4 & 12 (24) \\
CASCI(2,2) [classical] & $-77.116608$ & --- & --- \\
\bottomrule
\end{tabular}
\end{table}
\footnotetext{Hardware-efficient optimization reached the COBYLA iteration
cap (300/300); the reported energy should be treated as a lower-confidence
result pending a rerun with a larger budget.}

\subsection{Excited-State (Vertical Excitation) Energies}
Table~\ref{tab:excited} reports the QEOM excited states at the planar
geometry. Energy ordering and magnitude are consistent with a spin-resolved
CASCI(2,2) calculation performed as a cross-check (triplet at 4.91~eV,
singlet V-state at 14.24~eV, confirmed via $\langle S^2\rangle$
expectation values), allowing unambiguous spin assignment of the QEOM roots.

\begin{table}[h]
\centering
\caption{QEOM excited states, planar ethylene.}
\label{tab:excited}
\begin{tabular}{lccl}
\toprule
State & Energy (Ha) & $\Delta E$ (eV) & Assignment \\
\midrule
Excited 1 & $-76.941013$ & 4.779 & Triplet (T$_1$) \\
Excited 2 & $-76.596800$ & 14.145 & Singlet V-state ($\pi\to\pi^*$) \\
Excited 3 & $-76.407974$ & 19.283 & Doubly-excited singlet \\
\bottomrule
\end{tabular}
\end{table}

The triplet energy (4.78~eV) agrees well with the experimental vertical
triplet energy of ethylene ($\sim$4.6--4.9~eV). The singlet V-state
(14.1~eV) substantially overestimates the experimental value of 7.6~eV; this
is a known limitation of the minimal STO-3G basis combined with a (2,2)
active space for this particular ionic transition, and is not attributable
to the quantum algorithm.

\subsection{Classical Cross-Validation (CASCI)}
Table~\ref{tab:casci} lists the CASCI(2,2) results used for cross-validation.

\begin{table}[h]
\centering
\caption{Classical CASCI(2,2)/STO-3G benchmark.}
\label{tab:casci}
\begin{tabular}{lcc}
\toprule
Geometry & $E_0$ (Ha) & $\Delta E$, first excited root (eV) \\
\midrule
Planar (0\textdegree) & $-77.116608$ & 4.91 \\
``Twisted'' (90\textdegree)$^{*}$ & $-74.062176$ & 1.91 \\
\bottomrule
\end{tabular}
\end{table}

\noindent $^{*}$See Limitations (Sec.~\ref{sec:limitations}): the twisted
geometry as constructed corresponds to a distorted structure rather than an
ideal 90\textdegree\ rotation, and these values should be read qualitatively,
not as a literal rotational barrier.

\subsection{Noise and Error-Mitigation Analysis}

\subsubsection{Real-Hardware Ground-State Study (ibm\_marrakesh)}
Table~\ref{tab:hw} shows the converged UCCSD ground state evaluated on real
IBM Quantum hardware at increasing resilience levels.

\begin{table}[h]
\centering
\caption{Ground-state energy on \texttt{ibm\_marrakesh} at increasing
resilience levels.}
\label{tab:hw}
\begin{tabular}{lcc}
\toprule
Method & Energy (Ha) & Error vs.\ exact (mHa) \\
\midrule
Exact (statevector) & $-77.116628$ & 0.000 \\
Raw (resilience level 0) & $-76.985608$ & 131.02 \\
Readout-mitigated, TREX (level 1) & $-77.004366$ & 112.26 \\
ZNE + TREX + twirling (level 2) & $-77.070221$ & 46.41 \\
\bottomrule
\end{tabular}
\end{table}

Error mitigation reduced the deviation from the exact energy by $\sim$2.8$\times$
(131~mHa $\to$ 46~mHa). The raw hardware error alone (131~mHa $\approx 3.6$~eV)
exceeds the entire triplet excitation energy, underscoring that mitigation is
not optional for chemically meaningful results on current devices.

\subsubsection{Local Noise Study of Excited States (QEOM)}
Table~\ref{tab:qeom_noise} shows the QEOM excitation energies recomputed
locally under shot noise only and under a realistic device noise model
(\texttt{FakeSherbrooke} snapshot), with the ground state held fixed at its
noiselessly-optimized value.

\begin{table}[h]
\centering
\caption{QEOM excitation energies under local noise conditions.}
\label{tab:qeom_noise}
\begin{tabular}{lcc}
\toprule
Condition & Triplet (eV) & Singlet V (eV) \\
\midrule
Exact & 4.779 & 14.145 \\
Shot noise only & 4.773 & 14.022 \\
Realistic device noise & 4.817 & 14.146 \\
\bottomrule
\end{tabular}
\end{table}

The excited-state drift here ($\leq$0.15~eV) is much smaller than the raw
hardware ground-state error (131~mHa $\approx 3.6$~eV), most likely because
the local model was transpiled against an all-to-all-connected simulator and
omits the SWAP-gate overhead that real device connectivity imposes. This
suggests routing overhead, not just intrinsic gate/readout error, dominates
the real-hardware error and merits explicit modeling in future work.

\section{Limitations}
\label{sec:limitations}
\begin{itemize}
\item \textbf{Twisted geometry:} the rotation-axis parametrization used for
the twisted structure distorts C--H bond lengths rather than performing a
pure rotation about the C--C axis, so the reported ``twisted'' CASCI energy
reflects a distorted structure, not an ideal 90\textdegree\ ethylene, and
should be read qualitatively only.
\item \textbf{Spin labeling:} QEOM roots are assigned triplet/singlet
character by cross-validation against a separate $\langle S^2\rangle$-tagged
CASCI run, not by spin-purification within QEOM itself.
\item \textbf{Basis set:} STO-3G with a (2,2) active space is known to
overestimate ionic $\pi\to\pi^*$ energies; a larger basis/active space
(e.g.\ 2$e$/4$o$) would likely narrow the 14.1 vs.\ 7.6~eV gap.
\item \textbf{Noise-model fidelity:} the local QEOM noise study omits
target-backend connectivity, understating real device error.
\item \textbf{Convergence:} the hardware-efficient ansatz hit its COBYLA
iteration cap without an explicit convergence flag.
\end{itemize}

\section{Extension: Toward Photo-Induced Charge Transfer in Larger Chromophores}
Extending this pipeline toward a photosynthesis-relevant charge-transfer
estimate (e.g.\ a chlorophyll/porphyrin fragment) requires three steps beyond
what is implemented here. \textbf{(1) Active-space reduction via embedding:}
a full porphyrin macrocycle is too large for near-term treatment, so we
propose selecting a minimal donor--bridge--acceptor active space (frontier
orbitals of a donor and acceptor site, analogous to ethylene's $\pi/\pi^*$
pair) via DMET or CASSCF-based selection, with the remaining macrocycle and
protein/solvent environment treated at a mean-field/QM/MM level.
\textbf{(2) Electronic coupling via VQE+qEOM:} the same machinery used here
can, with minor extension, obtain the coupling $H_{DA}$ between
donor-localized and acceptor-localized diabatic states, e.g.\ via a
Hadamard-test-style circuit or a generalized qEOM subspace spanning both
localized configurations. \textbf{(3) Rate theory and photocurrent:} with
$H_{DA}$, a driving force $\Delta G$ (from energy differences at donor/acceptor
geometries, as we did for planar/twisted ethylene), and a reorganization
energy $\lambda$ (from the same pipeline at relaxed vs.\ vertical geometries),
Marcus theory gives the transfer rate
\[
k_{ET} = \frac{2\pi}{\hbar}\,|H_{DA}|^2\,
\frac{1}{\sqrt{4\pi\lambda k_B T}}\,
\exp\!\left[-\frac{(\Delta G + \lambda)^2}{4\lambda k_B T}\right],
\]
and chaining $k_{ET}$ across a multi-site pigment network via a Lindblad
master equation yields a steady-state flux that, scaled by charge and
quantum yield, estimates a photocurrent. \textbf{Bonus:} a toy demonstration
on butadiene intermediate in size between ethylene and a porphyrin fragment is a natural, tractable next step before full-scale embedding.

\section{Conclusion}
We implemented and validated a complete VQE+QSE (qEOM) pipeline for ethylene,
cross-checked against classical CASCI, and characterized real-hardware and
simulated noise alongside standard error-mitigation techniques. The
ground-state result is essentially exact; the excited-state result correctly
recovers the triplet energy while exposing a well-understood basis-set
limitation for the singlet V-state. Error mitigation on real IBM hardware
reduced the ground-state error by $\sim$2.8$\times$, and we outlined a
concrete path from this pipeline toward electron-transfer rates in larger,
biologically relevant chromophore systems.

\end{document}